% APÊNDICE

\appendix
\chapter{Integração Numérica pelo Método de Newmark e Convergência por Newton-Raphson}
\label{ap:newmark_nr}

O Método de Newmark \cite{newmark} é um procedimento numérico robusto para solucionar equações diferenciais no domínio do tempo. Na análise de sistemas rotativos e engrenagens com severas não linearidades, sua formulação clássica é acoplada a rotinas iterativas de correção para garantir convergência e estabilidade da solução \cite{lais2022}.

\section{Formulação Clássica do Método}

A formulação baseia-se na expansão em série de Taylor para o vetor de deslocamentos $\{q\}$. Considerando o passo de tempo $\Delta t = t_2 - t_1$, o método introduz os parâmetros $\beta$ e $\gamma$ para controlar a precisão e estabilidade, relacionando os estados do sistema entre dois instantes consecutivos:
\begin{gather}
	\{q\}_{t_2} = \{q\}_{t_1} + \{\dot{q}\}_{t_1}\Delta t + \left( \frac{1}{2} - \beta \right) \Delta t^2 \{\ddot{q}\}_{t_1} + \beta \Delta t^2 \{\ddot{q}\}_{t_2}
\\
	\{\dot{q}\}_{t_2} = \{\dot{q}\}_{t_1} + (1 - \gamma)\Delta t \{\ddot{q}\}_{t_1} + \gamma \Delta t \{\ddot{q}\}_{t_2}
\end{gather}

Adotando a regra do ponto médio ($\beta = 1/4$ e $\gamma = 1/2$), garante-se um esquema incondicionalmente estável para sistemas lineares e precisão de segunda ordem \cite{lais2022}. O método é implícito, pois as equações dependem da aceleração no próprio instante a ser avaliado ($\{\ddot{q}\}_{t_2}$).

\newpage
\section{Formulação de Convergência e Newton-Raphson}

Devido às não linearidades a solução implícita direta não é possível, adotando-se uma abordagem preditor-corretor com Newton-Raphson. Na etapa de previsão, as estimativas para o próximo passo assumem um incremento nulo de aceleração:
\begin{gather}
	\{q\}^*_{t_2} = \{q\}_{t_1} + \{\dot{q}\}_{t_1}\Delta t + \left( \frac{1}{2} - \beta \right) \Delta t^2 \{\ddot{q}\}_{t_1}
\\
	\{\dot{q}\}^*_{t_2} = \{\dot{q}\}_{t_1} + (1 - \gamma)\Delta t \{\ddot{q}\}_{t_1}
\\
	\{\ddot{q}\}^*_{t_2} = \{0\}
\end{gather}

Em seguida, avalia-se o resíduo dinâmico $\{R\}$ do sistema. Para satisfazer o equilíbrio dinâmico, o resíduo deve tender a zero:
\begin{equation}
	\{R\}_{t_2} = \{F_{ext}\}_{t_2} + \{F_{nl}(\{q\}_{t_2}, \{\dot{q}\}_{t_2})\} - \left( [M]\{\ddot{q}\}_{t_2} + [C]\{\dot{q}\}_{t_2} + [K]\{q\}_{t_2} \right) = 0
\end{equation}

Sendo $[M]$, $[C]$ e $[K]$ as matrizes de massa, amortecimento e rigidez; $\{F_{ext}\}$ as forças externas e $\{F_{nl}\}$ as forças não lineares. 

Para anular o resíduo iterativamente, aplica-se o método de Newton-Raphson. A correção dos estados exige a inversão da matriz Jacobiana $[J]$ que, devido à complexidade analítica, é construída numericamente \cite{jacobiano}:
\begin{equation}
	[J] = [M] + \gamma \Delta t [C] + \beta \Delta t^2 [K] - \frac{\partial \{F_{nl}\}}{\partial \{\ddot{q}\}}
\end{equation}

A derivada numérica é obtida aplicando uma pequena perturbação $\epsilon$ nas acelerações e avaliando a variação da força não linear correspondente. A correção da aceleração ($\Delta \{\ddot{q}\}$) a cada iteração de Newton resolve o sistema:
\begin{equation}
	[J] \Delta \{\ddot{q}\} = \{R\}
\end{equation}

Os vetores de estado são então atualizados:
\begin{gather}
	\{\ddot{q}\}_{t_2} \leftarrow \{\ddot{q}\}_{t_2} + \Delta \{\ddot{q}\}
\\
	\{\dot{q}\}_{t_2} \leftarrow \{\dot{q}\}_{t_2} + \gamma \Delta t \Delta \{\ddot{q}\}
\\
	\{q\}_{t_2} \leftarrow \{q\}_{t_2} + \beta \Delta t^2 \Delta \{\ddot{q}\}
\end{gather}

O processo iterativo se repete até a norma do resíduo ser inferior à tolerância definida ($tol = 10^{-6}$).

\section{Algoritmo de Passo de Tempo Adaptativo}

Transições abruptas da força exigem controle do incremento de tempo. Para otimizar o custo computacional mantendo a estabilidade, utiliza-se um passo de tempo adaptativo baseado nas iterações do Newton-Raphson:

\begin{itemize}
	\item Convergência Rápida: Se atingida em poucas iterações (menor ou igual a $5$), o subpasso ($\Delta t_{sub}$) é duplicado, respeitando o limite do passo macro.
	\item Convergência Lenta: Se demandar muitas iterações (maior ou igual a $10$), o subpasso é reduzido à metade, com um limite mínimo para evitar interrupções.
	\item Divergência: Caso não atinja a tolerância em 15 iterações, o passo é descartado, o subpasso cai para $25\%$ de seu valor e a integração retrocede.
\end{itemize}

A rotina foi implementada em Python e compilada \textit{just-in-time} (JIT) via \texttt{Numba}, garantindo desempenho computacional.